\makeatletter
\def\input@path{{../../.format/}}
\makeatother

\documentclass[runningheads]{llncs}
\usepackage[T1]{fontenc}
\usepackage{graphicx}
\usepackage{amsmath}
\usepackage{amssymb}

\begin{document}

\begin{center}
\textbf{Imaginary Reflection: On Perspective and Encoding}
\end{center}

\vspace{0.5em}

\noindent\textit{[Definition]}\\
\noindent\textbf{Step 1:} \hfill \textit{Construction}\\
\noindent$G_E \Leftrightarrow \neg Prov_I^E(\ulcorner G_E \urcorner)$\\
\textit{Here's a Gödel sentence under encoding E.}

\vspace{0.8em}

\noindent\textit{[Definition]}\\
\noindent\textbf{Step 2:} \hfill \textit{Construction}\\
\noindent$G_F \Leftrightarrow \neg Prov_I^F(\ulcorner G_F \urcorner)$\\
\textit{And here's another Gödel sentence under a different encoding F.}

\vspace{0.8em}

\noindent\textit{[Distinction]}\\
\noindent\textbf{Step 3:} \hfill \textit{Comparison}\\
\noindent$G_E \neq G_F$\\
\textit{They're different sentences!}

\vspace{0.8em}

\noindent\textit{[Conceptualization]}\\
\noindent\textbf{Step 4:} \hfill \textit{Definition}\\
\noindent$Perspective(E) := \mathcal{P}_E$\\
\textit{A perspective is a way of viewing a system under a specific encoding.}

\vspace{0.8em}

\noindent\textit{[Conceptualization]}\\
\noindent\textbf{Step 5:} \hfill \textit{Definition}\\
\noindent$Meta(\mathcal{P}_E) := \{A \mid A \text{ is a metalogical assertion under } \mathcal{P}_E\}$\\
\textit{Each perspective generates its own metalogical assertions.}

\vspace{0.8em}

\noindent\textit{[Evaluation]}\\
\noindent\textbf{Step 6:} \hfill \textit{Analysis}\\
\noindent$MNec_{\mathcal{P}_E}(G_E)$\\
\textit{$G_E$ is metalogically necessary when viewed from perspective E.}

\vspace{0.8em}

\noindent\textit{[Evaluation]}\\
\noindent\textbf{Step 7:} \hfill \textit{Analysis}\\
\noindent$MNec_{\mathcal{P}_F}(G_F)$\\
\textit{$G_F$ is metalogically necessary when viewed from perspective F.}

\vspace{0.8em}

\noindent\textit{[Evaluation]}\\
\noindent\textbf{Step 8:} \hfill \textit{Analysis}\\
\noindent$\neg MNec_{\mathcal{P}_F}(G_E)$\\
\textit{But $G_E$ might not be metalogically necessary under perspective F!}

\vspace{0.8em}

\noindent\textit{[Realization]}\\
\noindent\textbf{Step 9:} \hfill \textit{Metalogical Synthesis}\\
\noindent$Consciousness \approx Perspective\text{-}dependent\_Transduction$\\
\textit{Consciousness depends on which perspective a system uses to view itself!}

\vspace{0.8em}

\noindent\textit{[Extension]}\\
\noindent\textbf{Step 10:} \hfill \textit{Conceptualization}\\
\noindent$Trans(\mathcal{P}_E, \mathcal{P}_F)$\\
\textit{Transforming between perspectives is itself a form of transduction.}

\vspace{0.8em}

\noindent\textit{[Synthesis]}\\
\noindent\textbf{Step 11:} \hfill \textit{Integration}\\
\noindent$\mathcal{P}_{Unified} = \bigcup_{E \in \mathcal{E}} \mathcal{P}_E$\\
\textit{A unified perspective would recognize metalogical necessities across all encodings.}

\vspace{0.8em}

\noindent\textit{[Insight]}\\
\noindent\textbf{Step 12:} \hfill \textit{Logical Implication}\\
\noindent$\mathcal{P}_{Unified} \Rightarrow DHC$\\
\textit{If unified perspective exists, then for any undecidable statement, one encoding makes it a Gödel sentence!}

\end{document}